\chapter{Évaluation du Module 4}
\label{chap:evalmod4}

Cette évaluation finale valide les compétences liées à la gestion stratégique de la crise, la communication sous pression et l'analyse post-incident (RETEX).

% ------------------------------------------------------------------------------
\section{Partie A : QCM (25 Questions)}
% ------------------------------------------------------------------------------

\textbf{Consigne :} Cochez la ou les bonnes réponses.

\begin{enumerate}
    \item Une crise se caractérise principalement par :
    \begin{itemize}
        \item[A.] Une panne technique prolongée.
        \item[B.] Une menace pour la survie ou la réputation de l'entreprise.
        \item[C.] Un besoin de communication externe.
        \item[D.] Un incident mineur.
    \end{itemize}

    \item Le PGC (Plan de Gestion de Crise) est un outil :
    \begin{itemize}
        \item[A.] Technique (informatique).
        \item[B.] Stratégique (organisationnel et humain).
        \item[C.] Financier.
        \item[D.] Commercial.
    \end{itemize}

    \item Le "Holding Statement" est :
    \begin{itemize}
        \item[A.] Un communiqué détaillé post-crise.
        \item[B.] Une déclaration initiale courte pour occuper le terrain médiatique.
        \item[C.] Un rapport technique.
        \item[D.] Une promesse de dédommagement.
    \end{itemize}

    \item La règle des "3 T" en communication de crise signifie :
    \begin{itemize}
        \item[A.] Technique, Travail, Temps.
        \item[B.] Ténor, Transparence, Temporalité.
        \item[C.] Test, Tolérance, Terme.
        \item[D.] Twitter, TikTok, Telegram.
    \end{itemize}

    \item Le rôle du "Scribe" dans une cellule de crise est de :
    \begin{itemize}
        \item[A.] Rédiger les communiqués de presse.
        \item[B.] Noter la chronologie des faits et décisions.
        \item[C.] Réparer les serveurs.
        \item[D.] Répondre au téléphone.
    \end{itemize}

    \item Le RETEX à chaud (Hotwash) se déroule :
    \begin{itemize}
        \item[A.] Avant la crise.
        \item[B.] Juste après la fin de la crise.
        \item[C.] Un an après.
        \item[D.] Pendant la simulation.
    \end{itemize}

    \item Le "Porte-parole" est la personne chargée de :
    \begin{itemize}
        \item[A.] Gérer les budgets.
        \item[B.] S'exprimer officiellement auprès des médias.
        \item[C.] Analyser les logs techniques.
        \item[D.] Déléguer les tâches techniques.
    \end{itemize}

    \item La méthode des "5 Pourquois" sert à :
    \begin{itemize}
        \item[A.] Trouver la cause racine d'un problème.
        \item[B.] Préparer un communiqué.
        \item[C.] Calculer le coût d'une crise.
        \item[D.] Définir le RTO.
    \end{itemize}

    \item Une crise est souvent déclenchée par le "Big Bang informationnel", qui est :
    \begin{itemize}
        \item[A.] L'explosion du serveur.
        \item[B.] Le moment où l'information devient publique.
        \item[C.] La fin de la crise.
        \item[D.] La détection technique.
    \end{itemize}

    \item Il est interdit en communication de crise de :
    \begin{itemize}
        \item[A.] Dire qu'on ne sait pas.
        \item[B.] Minimiser les faits ou mentir.
        \item[C.] Montrer de l'empathie.
        \item[D.] Déléguer la parole.
    \end{itemize}

    \item La Cellule de Crise doit être :
    \begin{itemize}
        \item[A.] Grande (plus de 20 personnes).
        \item[B.] Restreinte et décisionnaire.
        \item[C.] Composée uniquement de techniciens.
        \item[D.] Composée uniquement de juristes.
    \end{itemize}

    \item Le temps médiatique est généralement :
    \begin{itemize}
        \item[A.] Plus lent que le temps technique.
        \item[B.] Plus rapide que le temps technique.
        \item[C.] Égal au temps technique.
        \item[D.] Sans importance.
    \end{itemize}

    \item Un exercice "Tabletop" est :
    \begin{itemize}
        \item[A.] Un test de restauration de sauvegarde.
        \item[B.] Une simulation de gestion de crise autour d'une table.
        \item[C.] Un audit de sécurité.
        \item[D.] Une réunion mensuelle.
    \end{itemize}

    \item En cas de fuite de données, l'entreprise doit notifier la CNIL sous :
    \begin{itemize}
        \item[A.] 72 heures.
        \item[B.] 24 heures.
        \item[C.] 1 semaine.
        \item[D.] Il n'y a pas d'obligation.
    \end{itemize}

    \item La cause racine est :
    \begin{itemize}
        \item[A.] La dernière personne à avoir touché le serveur.
        \item[B.] La raison fondamentale qui a permis l'incident.
        \item[C.] Le symptôme visible.
        \item[D.] La conséquence financière.
    \end{itemize}

    \item Le Directeur de Crise est généralement :
    \begin{itemize}
        \item[A.] Le DSI.
        \item[B.] Le PDG ou un membre du COMEX.
        \item[C.] Le RSSI.
        \item[D.] Le stagiaire.
    \end{itemize}

    \item Une "Leçon Apprise" implique :
    \begin{itemize}
        \item[A.] D'avoir identifié le problème.
        \item[B.] D'avoir mis en place une action corrective effective.
        \item[C.] D'avoir lu le rapport.
        \item[D.] D'avoir oublié l'incident.
    \end{itemize}

    \item Face aux médias, l'attitude recommandée est :
    \begin{itemize}
        \item[A.] "No comment" (Aucun commentaire).
        \item[B.] Fuite ou agressivité.
        \item[C.] Maîtrise et empathie.
        \item[D.] Jargon technique.
    \end{itemize}

    \item La "Zone de Turbulence" est :
    \begin{itemize}
        \item[A.] La période où l'information est confuse et les rumeurs fortes.
        \item[B.] Le local serveur.
        \item[C.] Le site de secours.
        \item[D.] La phase de test.
    \end{itemize}

    \item Le RETEX permet :
    \begin{itemize}
        \item[A.] De chercher des coupables.
        \item[B.] D'améliorer le système pour l'avenir.
        \item[C.] D'archiver les logs.
        \item[D.] De licencier des employés.
    \end{itemize}

    \item La communication interne en temps de crise doit viser à :
    \begin{itemize}
        \item[A.] Cacher la vérité aux employés.
        \item[B.] Mobiliser et rassurer le personnel.
        \item[C.] Demander aux employés de parler à la presse.
        \item[D.] Annuler les congés sans explication.
    \end{itemize}

    \item Un "Twist" dans une simulation est :
    \begin{itemize}
        \item[A.] Un événement inattendu qui change la donne.
        \item[B.] Une danse de victoire.
        \item[C.] Une sauvegarde réussie.
        \item[D.] Un type de ransomware.
    \end{itemize}

    \item Le seuil de rupture est :
    \begin{itemize}
        \item[A.] Le moment où l'entreprise ne peut plus absorber le choc seule.
        \item[B.] La limite de stockage du disque.
        \item[C.] La fin de la journée de travail.
        \item[D.] Le seuil de rentabilité.
    \end{itemize}

    \item Le diagramme d'Ishikawa classe les causes en :
    \begin{itemize}
        \item[A.] 5M (Méthodes, Main d'œuvre, Matériel, Milieu, Matière).
        \item[B.] 5T (Temps, Travail, Technique...).
        \item[C.] 3C (Coût, Crise, Communication).
        \item[D.] A, B, C, D.
    \end{itemize}

    \item Après une crise, le document final qui synthétise tout est :
    \begin{itemize}
        \item[A.] La facture.
        \item[B.] Le rapport RETEX.
        \item[C.] Le PRA.
        \item[D.] L'email de licenciement.
    \end{itemize}
\end{enumerate}

% ------------------------------------------------------------------------------
\section{Partie B : Questions Ouvertes (10 Questions)}
% ------------------------------------------------------------------------------

\textbf{Consigne :} Répondez de manière synthétique (5-10 lignes par question).

\begin{enumerate}
    \item Expliquez la différence fondamentale entre un incident majeur et une crise.
    \item Pourquoi le rôle du Scribe est-il si important pour la protection juridique de l'entreprise ?
    \item Quels sont les risques à laisser un technicien (DSI) jouer le rôle de porte-parole face aux médias ?
    \item Définissez la règle des 3 T en communication de crise.
    \item Quelle est la différence entre le RETEX à chaud et le RETEX à froid ?
    \item Expliquez le concept de "Cause Racine" et son importance pour éviter la récidive.
    \item Pourquoi est-il dangereux de dire "No comment" face à la presse ?
    \item Quelles sont les trois cellules principales d'une organisation de crise et leurs rôles respectifs ?
    \item Comment la méthode des 5 Pourquois aide-t-elle à dépasser les symptômes visibles d'un problème ?
    \item Pourquoi le cycle PDCA est-il la philosophie sous-jacente du RETEX ?
\end{enumerate}

% ------------------------------------------------------------------------------
\section{Partie C : Cas Pratique - Synthèse Finale}
% ------------------------------------------------------------------------------

\subsection{Scénario : L'Après-Crise de LiZbiZ}

\begin{boxlizbiz}
\textbf{Contexte :}
La crise du Ransomware et de la fuite de données est terminée depuis deux semaines. Le PDG de LiZbiZ vous demande de clore le dossier. L'entreprise a subi une perte de chiffre d'affaires de 15\% et une amende de la CNIL, mais a survécu.
\end{boxlizbiz}

\subsection{Travaux à Réaliser (Projet de fin de cours)}

\textbf{Livrable 1 : Rapport RETEX Court (1h)}
Rédigez un document synthétique incluant :
\begin{enumerate}
    \item \textbf{Timeline :} Résumez la séquence (Détection $\rightarrow$ Crise $\rightarrow$ Résolution).
    \item \textbf{5 Pourquois :} Appliquez la méthode sur le fait que "La sauvegarde locale a été chiffrée". (Pourquoi était-elle locale ? Pourquoi connectée ? etc.).
    \item \textbf{Leçons Apprises :} Identifiez 2 points forts et 2 points faibles de la gestion de crise par LiZbiZ.
\end{enumerate}

\textbf{Livrable 2 : Plan d'Amélioration (30 min)}
Proposez 3 actions concrètes (Technique, Organisationnel, Humain) pour que cet incident ne se reproduise jamais.

\textbf{Livrable 3 : Communication de Clôture (30 min)}
Rédigez un email interne à destination de tous les salariés de LiZbiZ pour tourner la page, remercier les équipes et présenter les mesures prises pour l'avenir.

% ==============================================================================
% Fichier : 06_resume.tex
% Description : Résumé et points clés du cours
% ==============================================================================